\paragraph{Cayley--Poisson 乘极展开与 \texorpdfstring{$\chi^2$}{chi2} 能量谱}
以下结果在任意有限均值的输入测度下给出 \(\tilde h_t/\tilde g_t\) 的精确傅里叶展开、\(\chi^2\) 与 Rényi 阶 \(2\) 的闭式能量谱，以及乘极矩序列的尺度半群与耗散恒等式。

\begin{proposition}[Cayley--Poisson 乘极展开与傅里叶矩系数]\label{prop:xi-poisson-cauchy-cayley-multipole-fourier}
设 \(\tilde\nu\in\mathrm{Prob}(\RR)\) 具有有限均值
\[
\bar\gamma:=\int_{\RR}\gamma\,\dd\tilde\nu(\gamma).
\]
对 \(t>0\) 定义 Cauchy 核与混合密度
\[
\mathsf P_t(x):=\frac{1}{\pi}\frac{t}{x^2+t^2},
\qquad
\tilde g_t(x):=\mathsf P_t(x-\bar\gamma),
\qquad
\tilde h_t(x):=(\mathsf P_t*\tilde\nu)(x)=\int_{\RR}\mathsf P_t(x-\gamma)\,\dd\tilde\nu(\gamma),
\]
并记密度比 \(u_t(x):=\tilde h_t(x)/\tilde g_t(x)\)。
定义 Cayley 相位变量
\[
\zeta_t(x):=\frac{(x-\bar\gamma)-\mathrm{i}t}{(x-\bar\gamma)+\mathrm{i}t}\in\TT,
\]
以及乘极变量
\[
W_t(\gamma):=\frac{\gamma-\bar\gamma}{(\gamma-\bar\gamma)+2\mathrm{i}t}\in\mathbb D.
\]
则对 \(\zeta=\zeta_t(x)\in\TT\) 成立严格恒等式
\[
u_t(x)
=\EE\!\left[\frac{1-|W_t(\Gamma)|^2}{|\zeta-W_t(\Gamma)|^2}\right],
\qquad \Gamma\sim\tilde\nu,
\]
并且 \(u_t\) 在 \(\TT\) 上具有全傅里叶展开
\[
u_t(\zeta)=1+\sum_{k\ge1}\Bigl(m_k(t)\,\zeta^{-k}+\overline{m_k(t)}\,\zeta^{k}\Bigr),
\qquad
m_k(t):=\EE\!\left[W_t(\Gamma)^k\right].
\]
其中 \(m_k(t)\) 对所有 \(k\ge1\) 均无需高阶矩条件即可定义。
\end{proposition}
\begin{proof}
令 \(y:=x-\bar\gamma\)、\(a:=\gamma-\bar\gamma\)。
直接化简得到
\[
\frac{\mathsf P_t(y-a)}{\mathsf P_t(y)}=\frac{y^2+t^2}{(y-a)^2+t^2}.
\]
另一方面取 \(\zeta=\frac{y-\mathrm{i}t}{y+\mathrm{i}t}\)、\(w=\frac{a}{a+2\mathrm{i}t}\)。
则
\[
1-|w|^2=\frac{4t^2}{a^2+4t^2},
\qquad
|\zeta-w|^2=\frac{4t^2\bigl((y-a)^2+t^2\bigr)}{(y^2+t^2)(a^2+4t^2)}.
\]
故
\[
\frac{1-|w|^2}{|\zeta-w|^2}=\frac{y^2+t^2}{(y-a)^2+t^2}=\frac{\mathsf P_t(y-a)}{\mathsf P_t(y)}.
\]
对 \(a=\Gamma-\bar\gamma\) 取期望即得第一式。
\par
对 \(|w|<1\)、\(|\zeta|=1\) 的 Poisson 核标准傅里叶级数为
\[
\frac{1-|w|^2}{|\zeta-w|^2}
=1+\sum_{k\ge1}\bigl(w^k\zeta^{-k}+\overline{w}^{\,k}\zeta^{k}\bigr),
\]
而 \(|W_t(\gamma)|<1\) 对所有实 \(\gamma\) 成立，故逐项取期望合法并得到所示展开。
\end{proof}

\begin{proposition}[\texorpdfstring{$\chi^2$}{chi2} 散度的谱恒等式]\label{prop:xi-poisson-cauchy-chi2-spectrum-multipole}
在命题 \ref{prop:xi-poisson-cauchy-cayley-multipole-fourier} 的设定下，令
\[
\chi^2(\tilde h_t\mid\tilde g_t)
:=\int_{\RR}\frac{(\tilde h_t-\tilde g_t)^2}{\tilde g_t}\,\dd x
=\int_{\RR}\tilde g_t(x)\,(u_t(x)-1)^2\,\dd x.
\]
则有严格恒等式
\[
\boxed{\;
\chi^2(\tilde h_t\mid\tilde g_t)=2\sum_{k\ge1}|m_k(t)|^2\; }.
\]
\end{proposition}
\begin{proof}
由 Cayley 变量 \(\zeta_t\) 的 Haar 回拉恒等式（命题 \ref{prop:app-busemann-poisson-minus1} 的等价表述），测度 \(\tilde g_t(x)\,\dd x\) 推前为 \(\TT\) 上 Haar 测度 \(\dd\theta/(2\pi)\)。
因此
\[
\chi^2(\tilde h_t\mid\tilde g_t)=\int_{0}^{2\pi}\bigl(u_t(e^{\mathrm{i}\theta})-1\bigr)^2\,\frac{\dd\theta}{2\pi}.
\]
由命题 \ref{prop:xi-poisson-cauchy-cayley-multipole-fourier}，\(u_t-1\) 的傅里叶系数满足 \(c_{-k}=m_k(t)\)、\(c_k=\overline{m_k(t)}\)（\(k\ge1\)），而 \(c_0=0\)。
Parseval 恒等式给出
\[
\int_{0}^{2\pi}\bigl(u_t-1\bigr)^2\,\frac{\dd\theta}{2\pi}
=\sum_{k\in\ZZ}|c_k|^2
=2\sum_{k\ge1}|m_k(t)|^2.
\]
\end{proof}

\begin{corollary}[Rényi 阶 \(2\) 的闭式表示与对 KL 的非渐近上界]\label{cor:xi-poisson-cauchy-renyi2-kl-upper}
在上述设定下，Rényi 阶 \(2\) 相对熵满足闭式
\[
D_2(\tilde h_t\mid\tilde g_t)=\log\!\left(1+2\sum_{k\ge1}|m_k(t)|^2\right).
\]
并且由 \(D_{\mathrm{KL}}\le D_2\) 得到非渐近上界
\[
D_{\mathrm{KL}}(\tilde h_t\mid\tilde g_t)\le \log\!\left(1+2\sum_{k\ge1}|m_k(t)|^2\right).
\]
\end{corollary}
\begin{proof}
由定义
\[
D_2(\tilde h_t\mid\tilde g_t)=\log\int_{\RR}\frac{\tilde h_t(x)^2}{\tilde g_t(x)}\,\dd x
=\log\int_{\RR}\tilde g_t(x)\,u_t(x)^2\,\dd x.
\]
又因 \(\int \tilde g_t u_t=\int \tilde h_t=1\)，故
\[
\int \tilde g_t u_t^2
=\int \tilde g_t\bigl((u_t-1)+1\bigr)^2
=1+\int \tilde g_t(u_t-1)^2
=1+\chi^2(\tilde h_t\mid\tilde g_t).
\]
代入命题 \ref{prop:xi-poisson-cauchy-chi2-spectrum-multipole} 即得闭式。
最后使用 Rényi 相对熵关于阶的单调性（\(\alpha\ge1\) 时 \(D_1=D_{\mathrm{KL}}\le D_2\)）得到上界。
\end{proof}

\begin{proposition}[乘极矩的尺度演化链]\label{prop:xi-poisson-cauchy-multipole-moment-chain}
在命题 \ref{prop:xi-poisson-cauchy-cayley-multipole-fourier} 的设定下，对每个 \(k\ge1\)，函数 \(t\mapsto m_k(t)\) 在 \((0,\infty)\) 上可微，且满足精确微分链
\[
t\,\frac{\dd}{\dd t}m_k(t)=-k\bigl(m_k(t)-m_{k+1}(t)\bigr).
\]
\end{proposition}
\begin{proof}
固定 \(a\in\RR\)，令 \(w(t):=\frac{a}{a+2\mathrm{i}t}\)。
直接求导得
\[
w'(t)=-\frac{2\mathrm{i}a}{(a+2\mathrm{i}t)^2}.
\]
另一方面，
\(
w(t)\,(1-w(t))=\frac{2\mathrm{i}ta}{(a+2\mathrm{i}t)^2}
\)，故 \(t\,w'(t)=-w(t)\,(1-w(t))\)。
于是
\[
\frac{\dd}{\dd t}w(t)^k
=k\,w(t)^{k-1}w'(t)
=-\frac{k}{t}\bigl(w(t)^k-w(t)^{k+1}\bigr).
\]
取 \(a=\Gamma-\bar\gamma\) 并对 \(\Gamma\sim\tilde\nu\) 取期望；由于 \(|w(t)|<1\) 且导数有一致有界支配，可交换微分与期望，从而得到所示微分链。
\end{proof}

\begin{proposition}[\texorpdfstring{$\chi^2$}{chi2} 的精确耗散恒等式]\label{prop:xi-poisson-cauchy-chi2-dissipation-identity}
令 \(\chi^2(t):=\chi^2(\tilde h_t\mid\tilde g_t)\)。
则成立严格耗散恒等式
\[
\frac{\dd}{\dd t}\bigl(t\,\chi^2(t)\bigr)
=-2\sum_{k\ge1}k\,\bigl|m_k(t)-m_{k+1}(t)\bigr|^2\le 0.
\]
并且等号在某个 \(t>0\) 成立当且仅当 \(\tilde\nu=\delta_{\bar\gamma}\)（从而对所有 \(t>0\) 有 \(\tilde h_t\equiv \tilde g_t\)）。
\end{proposition}
\begin{proof}
由命题 \ref{prop:xi-poisson-cauchy-chi2-spectrum-multipole}，
\(
\chi^2(t)=2\sum_{k\ge1}|m_k(t)|^2
\)，从而
\[
\chi^2{}'(t)=4\sum_{k\ge1}\Re\bigl(m_k'(t)\overline{m_k(t)}\bigr).
\]
代入命题 \ref{prop:xi-poisson-cauchy-multipole-moment-chain} 的 \(m_k'(t)=-(k/t)(m_k-m_{k+1})\)，并整理得到
\[
\frac{\dd}{\dd t}\bigl(t\,\chi^2(t)\bigr)
=-2\sum_{k\ge1}k\,|m_k-m_{k+1}|^2.
\]
若右端在某个 \(t>0\) 为零，则 \(m_k(t)=m_{k+1}(t)\) 对所有 \(k\ge1\) 成立。
又由 \(|m_k(t)|\le \EE|W_t(\Gamma)|^k\to 0\)（\(k\to\infty\)）可得 \(m_k(t)=0\) 对所有 \(k\ge1\)。
由命题 \ref{prop:xi-poisson-cauchy-cayley-multipole-fourier} 的傅里叶展开得到 \(u_t\equiv 1\)，即 \(\tilde h_t\equiv \tilde g_t\)，从而 \(\tilde\nu=\delta_{\bar\gamma}\)。
反向蕴含显然。
\end{proof}

\begin{proposition}[尺度 Möbius 共变与 Koenigs 线性化]\label{prop:xi-poisson-cauchy-multipole-mobius-koenigs}
对任意 \(s>0\) 定义 Möbius 自映射
\[
\Phi_s(u):=\frac{u}{s-(s-1)u}\qquad(u\in\mathbb D).
\]
则对每个 \(\gamma\in\RR\) 与 \(t>0\) 有严格共变关系
\[
W_{st}(\gamma)=\Phi_s\bigl(W_t(\gamma)\bigr),
\]
并且 \((\Phi_s)_{s>0}\) 构成乘法参数化半群：\(\Phi_s\circ\Phi_r=\Phi_{sr}\)。
再者，Koenigs 坐标
\[
\kappa(u):=\frac{u}{1-u}
\]
将该半群对角化为纯伸缩：
\[
\kappa\!\left(W_t(\gamma)\right)=\frac{\gamma-\bar\gamma}{2\mathrm{i}t},
\qquad
\kappa(\Phi_s(u))=\frac{1}{s}\,\kappa(u).
\]
\end{proposition}
\begin{proof}
令 \(a:=\gamma-\bar\gamma\)，则 \(W_t(\gamma)=\frac{a}{a+2\mathrm{i}t}\)。
由此解出 \(a=2\mathrm{i}t\,\frac{W_t}{1-W_t}\)，从而
\[
W_{st}
=\frac{a}{a+2\mathrm{i}st}
=\frac{\frac{W_t}{1-W_t}}{\frac{W_t}{1-W_t}+s}
=\frac{W_t}{s-(s-1)W_t}
=\Phi_s(W_t).
\]
半群性质由直接代入或由 Koenigs 坐标下的线性化得到。
最后由 \(\kappa(W_t)=\frac{W_t}{1-W_t}=\frac{a}{2\mathrm{i}t}\) 得到对角化恒等式。
\end{proof}

\begin{proposition}[有限模态压缩的 \texorpdfstring{$L^2(\tilde g_t)$}{L2(gt)} 精确误差]\label{prop:xi-poisson-cauchy-multipole-truncation-l2}
对 \(N\ge1\) 定义截断
\[
u_{t,N}(\zeta):=1+\sum_{k=1}^{N}\Bigl(m_k(t)\,\zeta^{-k}+\overline{m_k(t)}\,\zeta^{k}\Bigr),
\qquad \zeta\in\TT,
\]
并令 \( \tilde h_{t,N}(x):=\tilde g_t(x)\,u_{t,N}\bigl(\zeta_t(x)\bigr)\)。
则有严格等式
\[
\int_{\RR}\frac{\bigl(\tilde h_t(x)-\tilde h_{t,N}(x)\bigr)^2}{\tilde g_t(x)}\,\dd x
=2\sum_{k>N}|m_k(t)|^2.
\]
从而 \((m_1(t),\dots,m_N(t))\) 构成 \(L^2(\tilde g_t)\) 意义下最优 \(N\)-模态近似的完备信息，残差能量完全由高频乘极能量尾项给出。
\end{proposition}
\begin{proof}
在 \(\TT\) 上以 Haar 内积，\(u_{t,N}\) 为 \(u_t\) 在次数不超过 \(N\) 的三角多项式子空间上的正交投影；
故
\[
\|u_t-u_{t,N}\|_{L^2(\TT)}^2=\sum_{|k|>N}|c_k|^2=2\sum_{k>N}|m_k(t)|^2.
\]
再由 \(\tilde g_t(x)\,\dd x\) 为 Haar 测度回拉得
\(
\|u_t-u_{t,N}\|_{L^2(\TT)}^2=\int(\tilde h_t-\tilde h_{t,N})^2/\tilde g_t
\)。
\end{proof}

\endinput

